% !TeX root = ../../Book2.tex
\section{双线性振荡积分}
\label{b2:sec:2503}

二维曲线的双线性估计可以用一个可逆的频率和映射完整计算。这给出横截增益的直接模型，并显示常数如何随方向靠近而恶化。

\subsection{双线性振荡积分}
\label{b2:subsec:250301}

令 \(\gamma(t)=(t,t^2)\)，
\[
 E_If(x,y)=\int_I e^{2\pi\mathrm i(xt+yt^2)}f(t)\,\mathrm dt.
\]
对两个不交紧区间 \(I,J\)，乘积的 Fourier 侧对应映射
\[
 \Psi(t,s)=(t+s,t^2+s^2),\qquad
 |\det D\Psi|=2|s-t|.
\]
由于 \(t\in I\)、\(s\in J\) 的顺序固定，和与平方和确定无序对 \(\{t,s\}\)，顺序又确定有序对，所以映射在该矩形内一对一。
这项一对一性不能只由 Jacobi 行列式非零推出；后者仅保证局部可逆。

\subsection{双线性估计}
\label{b2:subsec:250302}

\begin{theorem}[分离抛物线片的双线性估计]\label{b2:thm:25-parabola-bilinear}
若 \(\operatorname{dist}(I,J)\ge\rho>0\)，则
\[
 \|E_If\,E_Jg\|_{L^2(\mathbb R^2)}^2
 =\int_{I\times J}\frac{|f(t)|^2|g(s)|^2}{2|t-s|}\,\mathrm dt\,\mathrm ds
 \le(2\rho)^{-1}\|f\|_2^2\|g\|_2^2.
\]
\end{theorem}
\begin{proof}
先取光滑紧支撑输入。换元 \(u=\Psi(t,s)\) 后，乘积为某个密度 \(h(u)\) 的逆 Fourier 变换，且
\(h(\Psi(t,s))=f(t)g(s)/(2|t-s|)\)。
Plancherel 与换元给出
\(\|h\|_2^2=\int |f|^2|g|^2/(2|t-s|)\)，这就是等式。
右侧由分离下界控制。一般 \(L^2\) 输入以光滑函数逼近，两项延拓乘积在 \(L^2\) 中 Cauchy；同时每个延拓在点态积分意义下连续依赖于 \(L^1\) 输入，有限区间上的 \(L^2\) 逼近也给出 \(L^1\) 逼近，故极限确为目标乘积。
\end{proof}
这个定理估计乘积而非单个函数。它没有断言 \(E_I:L^2(I)\to L^4(\mathbb R^2)\) 在允许片自相互作用时成立；对角线 \(t=s\) 的 Jacobi 行列式退化正是该推断的障碍。

\subsection{曲率相互作用}
\label{b2:subsec:250303}

对一般曲线 \(\gamma\)，和映射 Jacobi 行列式为
\(\det(\gamma'(t),\gamma'(s))\)。在足够小且方向分离的两片上，若此 determinant 有正下界、映射具有受控有限重数，同一换元与 Plancherel 给出对应估计，常数还包含重数。
对于抛物线，曲率使不同参数的切向方向改变；对于两条本来就横截的直线，则不需要曲率。

\ProseParagraphBreak
因此双线性理论中的曲率与横截性各有角色。横截性直接保证频率和映射的 Jacobi 行列式，曲率决定同一曲面的不同片能否形成横截方向。高维中两个超曲面的和映射通常有正维纤维，需要沿纤维积分，不能照搬二维的一对一公式。
